% ==========================================================================
% Paper 2 — Section 9: The Unified Mechanistic Ladder & General Discussion.
% ==========================================================================
\section{The Unified Mechanistic Ladder and General Discussion}
\label{sec:ladder}

The eight preceding sections assembled one rung at a time. This section
assembles the whole: the master table of the ladder (§9.1), the two
orthogonal axes as the paper's conceptual residue (§9.2), the biological
reading of each organizational degree of freedom---explicitly as a mapping
of \emph{degrees of freedom}, not as an equivalence claim (§9.3)---and the
general discussion (§9.4).

\subsection{The master table}
\label{sec:master}

Table~\ref{tab:ladder} compresses the ladder. Each row names the single
added organizational degree of freedom, the primitive it certifies, the
boundary wall the frozen audits located at that rung, and the decisive
experiment. Two reading rules apply. First, the walls are results, not
footnotes: every $\not\Rightarrow$ in the table is backed by an exact
number (a zero flip rate, a positive softness distance, a positive binding
error). Second, the ladder is constructive in the tested family and
\emph{sufficient}, not necessary: the paper claims that these degrees of
freedom suffice here, never that they are the only possible ones.

\begin{table}[t]
\centering
\small
\setlength{\tabcolsep}{4pt}
\begin{tabular}{@{}p{1.4cm}p{2.7cm}p{3.0cm}p{3.2cm}p{3.2cm}@{}}
\toprule
Rung & Added DOF & Certified primitive & Boundary wall (frozen) &
Decisive audit \\
\midrule
L0 & (frozen baseline) & Memory & --- & Collision $K$: B1 $0.5$ exact vs B4
$2.047\times10^{4}$; post-reset $0$ vs $32.2\%$ \\
L1 & surprise-gated window & Saliency weighting & Order-locked:
$\operatorname{argmax}$ query-invariant & FlipRate $0.000$; Probe-2
\textsc{terminated} / \textsc{audit\_failed} \\
L2 & (measurement rung) & Geometric reorganization & Effective rank
$\neq$ intrinsic dimension; no selectivity & $d_D$ $2.089$ vs $1.484$;
shared-frame B4 $\not>$ B1 \\
L3 & O1 opponent competition & Nonmonotone selective tuning &
$\text{Competition}\not\Rightarrow\text{Routing}$ & Theorem 4.1; C1 peak
$x^*=1.00$; six layouts FlipRate $0.000$ \\
L4 & O2 vector Q-K & Content routing &
representation-family-dependent & Flip law $\theta/\pi$ (132/132);
$+0.2587\to-0.1559$; $0.4425$ vs $0.4445$; AM2 \\
L5 & O3 K-V + O4 WTA & Binding; exact delivery &
softmax mixing barrier; exactness only in the discrete limit &
$T_{\mathrm{swap}}=-T$ ($10^{-12}$); $\gamma$ ladder $D\to0$; flat
flip rate $0.4445$ \\
L6 & O5 parallel channels + combiner & Composition &
$\text{Binding}\not\Rightarrow\text{Composition}$; unconditional error
$=$ routing error & M0/M1/M2 chain; 8 counterfactuals;
$E_{\mathrm{comp}}|\text{both-routed}\equiv0$ \\
\bottomrule
\end{tabular}
\caption{The unified mechanistic ladder. Each row is one frozen sprint;
the walls are the paper's negative results, each with its exact
statistic.}
\label{tab:ladder}
\end{table}

\begin{figure}[t]
\centering
\includegraphics[width=0.72\textwidth]{fig6_ladder}
\caption{The unified mechanistic ladder of dynamical computation. Each
rung names the single added organizational degree of freedom and, in
italics, the boundary wall located at that rung by the frozen audits.}
\label{fig:ladder}
\end{figure}

\subsection{The two orthogonal axes}
\label{sec:axes}

The ladder's most transferable residue is a pair of orthogonal axes that
the audits forced into view.

\emph{Axis 1: routing selectivity $\perp$ readout discreteness} (§7).
Sharpening from soft mixing to one-hot selection changes \emph{how} a
winner is delivered, never \emph{which} winner is selected: the flip rate
is exactly flat ($0.444527$) across the entire WTA ladder, while the
softness distance collapses from $1.74$ to $0$. A model's reordering
capacity and its symbol-delivery capacity are independent purchases.

\emph{Axis 2: binding fidelity $\perp$ combiner error} (§8). The combiner
contributes exactly zero error conditioned on correct operand delivery
(measured maximum $0.0\mathrm{e}0$ over $3\times10^4$ histories), while
the unconditional composition error is entirely the propagation of
upstream routing errors. A model's ability to compose and its ability to
bind are independent purchases, separated by the operand-delivery step.

Together the axes state the paper's architectural thesis in measurable
form: capability in this substrate is organized along \emph{named,
independent} degrees of freedom, so that no single scalar (accuracy, loss,
rank, sharpness) can stand in for any of them---each rung owns its own
audit, and the audits are the ladder.

\enlargethispage{2\baselineskip}
\subsection{Biological reading: degrees of freedom, not equivalences}
\label{sec:bio}

The paper's organizational degrees of freedom have natural biological
counterparts, and the blueprint fixes the mapping to be read strictly as
an analogy between \emph{kinds of organization}, never as a claim that the
tested neuron is any particular biological circuit.

\begin{itemize}[leftmargin=*,noitemsep,topsep=1pt]
\item \emph{Opponent competition (L3).} The E-I microcircuit is the
canonical organizational motif of cortical columns
\citep{mountcastle1997}, where fast-spiking parvalbumin-positive
interneurons provide the subtractive inhibition whose functional
signature---nonmonotone, band-pass selectivity---Theorem 4.1 derives from
two inequalities alone \citep{amari1977}.
\item \emph{The two-dimensional query direction (L4).} A minimal
continuous direction space over candidate content is what oscillatory
phase (a circular variable) provides to relational codes
\citep{singer1999}; the paper's content is not that phase \emph{is}
$\mathbb{R}^2$, but that a connected one-dimensional query manifold is the
minimal geometry in which query-conditioned reordering becomes continuous.
\item \emph{Parallel binding channels (L5--L6).} Independent nonlinear
channels converging on a common readout are the organizational signature
of dendritic computation: apical and basal compartments of pyramidal
neurons integrate separate streams before somatic output
\citep{london2005,poirazi2003}. The paper's dual-channel-plus-combiner
architecture is the minimal abstraction of exactly that arrangement, and
the axon initial segment's threshold integration is the natural candidate
for the two-input algebraic node.
\end{itemize}

The caution inherited from the frozen discipline: these are candidate
physical realizations of non-factorizable organization, not established
identities, and none of the tested numbers transfers to biology by
construction (§8.6).

\subsection{General discussion}
\label{sec:discussion}

\emph{More dynamics does not imply more computational primitives.} The
temporal memory window $W=5$ was held strictly constant throughout the ladder:
each rung does not enlarge temporal buffer capacity, but systematically introduces
distinct \emph{organizational degrees of freedom}---a second sign of coupling
(E-I opponent units), a second query dimension ($S^1$ rotation), a split between
keys and values, a sharp selection map, a second channel, and one algebraic
combiner. These additions reflect how attention and memory structures function
as fast-weight programming and energy-based associative retrieval
\citep{schlag2021,ramsauer2020}. Every rung's effect is certified by a
counterfactual that the previous rung fails. Nonlinear dynamics is
the stage on which computation happens; the primitives are the topology
erected on it.

Three questions delimit the honest reach of the work. First,
\emph{necessity}: the paper establishes sufficiency chains inside one
frozen family; it does not rule out other mechanisms that route, bind, or
compose. Second, \emph{generality}: every boundary wall is certified at
frozen parameters, frozen queries, and two to four candidates; scaling
audits (more distractors, nested operators, learned embeddings) are
explicitly out of scope and listed in §10. Third, \emph{composition's
open end}: the combiner tested is a single two-input algebraic node; what
minimal organization suffices for multi-operand or nested composition is
the natural next rung, and the paper leaves it unclaimed.

What the paper does claim, and defends with exact numbers, is a boundary
map: for each rung from memory to composition, the minimal organizational
condition under which the primitive appears in this substrate, and the
exact statistic at which the previous rung stops. That map, not any single
model, is the intended contribution.
